\subsection{Negative drift on the primary scale}
\label{sec:direct-primary-negative}

Assume the magnitude of the constant coefficient has a positive
normalized lower bound on the primary scale. Positive boundary roots
may still merge. At bounded radii the drift-weighted curvature estimate
continues through this fold; at early heights and large exits we use
contraction instead. We prove that these ranges overlap and cover
\(\mathcal H_p\), so that \cref{cor:ordered-cover-zero-count} applies
to a single logarithmic slope ratio on the full interval.

\begin{proposition}[Primary negative drift]
\label{prop:direct-primary-negative}
For every fixed \(a_0^{\rm pr}>0\), the return displacement has at
most three distinct zeros when
\begin{equation}
\label{eq:direct-primary-region}
 q_0=-aH^3,\quad q_1=bH^2,\quad q_2=\bar c H,\qquad
 a_0^{\rm pr}\le a\le1,\quad 0\le b\le1,\quad |\bar c|\le1,
\end{equation}
for all sufficiently small \(H>0\) and original parameters.
\end{proposition}

Write \(|\eta_\sigma|=HX,\ z=H^2Y\). The outward
boundary drifts are
\begin{equation}
\label{eq:direct-primary-drifts}
 K_\sigma(X)=-H^{-3}q(\sigma HX)
 =a-\sigma bX-\bar c X^2-\sigma X^3R_q(\sigma HX).
\end{equation}
With the central-height scale \(Q=H^3\) of \cref{tab:direct-scale-guide},
the common logarithmic height coordinate below is \(T=-H\log(h/Q)\).
The coefficient calculation in the appendix shows that the two
drifts converge, on fixed boxes and in every required state order, to
\[
 K_+^0=a-bX-\bar c X^2+X^3,\qquad
 K_-^0=a+bX-\bar c X^2-X^3.
\]
For \(X>0\), their curvature coefficients are \(b/X^2+3\) and \(-b/X^2-3\),
respectively, so the endpoint estimates below must preserve these
opposite signs uniformly as positive boundary roots merge.
The drift \(K_-^0\) has a unique, uniformly simple positive root \(\beta\)
separated from zero. Its derivative has at most one positive
zero, and
\[
 -\frac{(K_-^0)'(\beta)}{\beta}
        =2a/\beta^2+b/\beta+\beta>0 .
\]
Compactness gives the uniform bounds, which persist under the
\(O(H)\) perturbation. All positive roots of \(K_+^0\), if
present, also lie in a fixed compact interval away from zero.

On every fixed positive radius interval, the height equation is
\[
 X_{\log Y}=HF_\sigma(X)+YB(\sigma HX,H^2Y)/X,\qquad
 F_\sigma=K_\sigma/X.
\]
Thus \cref{lem:direct-boundary-flow} applies with \(\delta=H\), and
\eqref{eq:direct-projection-energy} holds for its inverse projection.
Below the first root of \(F_\sigma\), where \(F_\sigma>0\),
\cref{rem:direct-projected-continuation} applies.

\subsubsection{Early central heights}

Choose a small fixed \(d_f>0\) and put \(z_e=H^3d_f\). The present
scales are \eqref{eq:direct-quadratic-prefix-scales} with
\(\delta=\ell=H\) and \(\bar c\) replaced by \(-\bar c\). By
\cref{lem:early-central-heights}, there is a small fixed \(R>0\) such
that the set \eqref{eq:direct-primary-early}, with \(r_0=R\), is an
initial interval in \(T\) on which the return contracts.

\subsubsection{Bounded radii, including the fold}

For every remaining central height, the positive half-orbit reaches
the radius \(R\) before \(z_e\). Fix a small \(Y_f>0\) for the central
entry and the negative-root box below, and let \(M\) be large; \(M\) is
fixed after \(M_1\) in the covering step. The same \(Y_f\) serves every
such \(M\): at \(H=0\) the positive projection on \([R/2,M]\) is
\(w^2=X^2-2B_0Y\), and only its lower-radius margin matters, while a
larger \(M\) requires only a smaller final \(H\) for the derivative
bounds.
Consider those central heights with \(X_+(H^2Y_f;h)\le M\).
Their projected state \(w_+\) at this height lies in \([R/2,M]\):
its value at \(z_e\), when used, differs from \(X_+(z_e)\)
by \(O(Hd_f)\), and \(w_+\) increases.
The phase comparison below also places their negative-branch state
\(w_-\) in \([d_0,\beta)\) for a fixed \(d_0>0\).

Choose a small fixed \(x_L>0\) inside the central positive-drift margins
of both \(K_\sigma\). With
\[
 \mathcal A_\sigma(x)=\int_0^x\frac{u\,du}{K_\sigma(u)},
\]
\eqref{eq:direct-common-phase} gives
\begin{equation}
\label{eq:direct-primary-phases}
 w_\sigma=\Phi_{F_\sigma}^{\,T+C_\sigma+
                              H\Xi_\sigma(\xi_\sigma)}(x_L),
 \quad
 C_\sigma=H\log(Y_f/H)-\mathcal A_\sigma(x_L),\quad
 \xi_\sigma=\e^{\mathcal A_\sigma(x_L)/H}h/H^3 .
\end{equation}
The phase corrections \(\Xi_\sigma\) vanish at zero and have
uniformly bounded derivatives of the required orders.
Positive advance to \(R/2\) takes a fixed positive boundary
time. Take both entry actions smaller than one quarter of
this minimum time. The two phases in
\eqref{eq:direct-primary-phases} then have a common positive
lower bound. The negative scalar flow advances into a fixed
interval away from zero and remains below its first root.
Thus \(d_0\le w_-<\beta\), and the inverse projection
continues the negative graph to \(Y_f\).

At \(t=H^2Y_f\), all three coefficient ratios in
\eqref{eq:direct-tail-smallness} are \(O(H)\), and
\cref{lem:boundary-flow-endpoint} with \(\ell=\delta=H\) gives
\begin{equation}
\label{eq:direct-primary-endpoints}
 -I_{\sigma,w}=E A_\sigma^*wV_\sigma(w),\qquad
 E=H^2(z_*/(H^2Y_f))^{-2q_2}\asymp H^2,\qquad
 A_\sigma^*\asymp1,\quad V_\sigma=1+O_{C^3}(H).
\end{equation}
The negative endpoint estimate must also hold in the large-exit range,
before the positive radius cut is chosen.
The positive states range over \([R/2,M]\), so the constants of the
positive endpoint depend on \(M\). The negative states lie in
\([d_0,\beta)\), and the constants of the negative endpoint are
independent of \(M\).

The drifts \eqref{eq:direct-primary-drifts} have
\(K_\sigma'(0)=-\sigma b\) and, in the notation of
\eqref{eq:direct-curvature-coefficient-cubic},
\(R_{K_\sigma}(w)=-\sigma R_q(\sigma Hw)\), which is
\(\sigma+O_{C^2}(H)\) on the state intervals by the coefficient
calculation in the appendix. Hence
\eqref{eq:direct-curvature-coefficient-cubic} gives
\begin{equation}
\label{eq:direct-primary-opposite-curvatures}
 \frac{wK_+''-K_+'}{w^2}=\frac b{w^2}+3+O_M(H)>2,\qquad
 \frac{wK_-''-K_-'}{w^2}=-\frac b{w^2}-3+O_M(H)<-2 .
\end{equation}
The margins hold also when \(b=0\). Since the negative state stays in
\([d_0,\beta)\), \cref{lem:curvature-criterion} gives \(\mathcal L''<0\)
in \(T\), uniformly up to either root. The bounded-radius central
heights form an interval, by \cref{lem:central-height-order}.

\subsubsection{Escaping radii and their endpoint slope}

Fix \(M_0\) beyond all positive roots and critical points of
the limiting primary polynomial. For the positive radius
\(\varrho=\eta\ge M_0H\), write
\[
 \varrho^+=H+\varrho,\qquad u=H/\varrho,\qquad
 v=\varrho/(H+\varrho).
\]
Then
\[
 -q(\varrho)=\varrho^2\varrho^+\mathcal K,\qquad
 \mathcal K=v\{au^3-bu^2-\bar c u-R_q(\varrho)\}.
\]
For large \(M_0\) this \(\mathcal K\) is positive and has
uniform derivative bounds on an enlarged base box.
It equals \(1+O(u+H+\varrho)\). Thus
\cref{lem:direct-adaptive-positive} applies with \(s=\varrho_0=H\).
Its exit and landing satisfy
\(\varrho_e\le\varrho_l\le2\varrho_e\), with uniform derivatives,
before the admissible endpoint.

If the thin exit is at \(\varrho_e\ge2M_0H\), there is no positive
boundary root before \(M_0H\). Indeed at a thin hit of
\(M_0H\) the primary height is \(Y=O_{M_0}(H)\). The
projection, constructed on the fixed interval up to \(M_0\),
differs from the radius by \(O_{M_0}(H)\) there.
Since any first root is separated from \(M_0\), a projected state below
that root cannot give such a hit.

On these no-root parameters, transport the regular central
height coordinate from \(x_L\) to \(M_0\), using
\eqref{eq:direct-bottleneck-action}. The endpoint collars
have fixed positive drift margins, although the action between
them may diverge. In the height \(y=z/H^3\), this gives
\begin{equation}
\label{eq:direct-primary-bottleneck-prefix}
 \Psi_o(y_o)=\e^{\mathcal A_+(M_0)/H}\frac h{H^3}
 \Upsilon\!\left(\e^{\mathcal A_+(x_L)/H}\frac h{H^3}\right),
 \qquad \Psi_o(y)=y+O(y^2),\quad \Upsilon(0)=1 .
\end{equation}
By \eqref{eq:direct-bottleneck-inverse-composition}, the units and
height coordinates, together with their inverses, have uniform
derivative bounds in the expanded input height.

There is a fixed positive boundary-action gap from \(M_0\)
to \(2M_0\). Hence an exit beyond \(2M_0H\) makes the
height at \(M_0H\), normalized by \(H^3\), at most
\(C H^{-N}\e^{-\mathcal A_{\rm gap}/H}\). The unscaled exit radius is
bounded, so its prefactor costs only a fixed power of \(1/H\).
Differentiating the regular units in
\eqref{eq:direct-primary-bottleneck-prefix} gives
\[
 \frac{d\log\widehat z(M_0H)}{d\log h}
             =1+O(\e^{-\mathcal A_{\rm gap}/(2H)}).
\]
This is the bound on \(d\lambda_{\rm in}/d\log h\) required in
\eqref{eq:direct-adaptive-source-slope}, with \(\varrho_i=M_0H\). The tail
condition there holds at the landing: \(\varrho_l\ge M_0H\) gives
\(|q_0|/(m_{\rm land}\varrho_l^2)+q_1/\sqrt{m_{\rm land}\varrho_l^2}=O(H)\), and
\(|q_2|\le H\).

To compare with the positive endpoint slope, return to the negative
branch. It retains a fixed positive entry-time gap; its projected state lies
in \([d_0,\beta)\), regardless of the large positive exit.
The independently chosen negative endpoint bounds in
\eqref{eq:direct-primary-endpoints} therefore give
\[
 0<m_-=hI_{-,h}\le CHE\le CH^3 .
\]
Thus \cref{cor:large-exit-contraction} applies, with \(s=H\), \(C_2=1\)
and \(C_1\) the constant just obtained, to the central heights in
\(\cH_p\) with \(\varrho_e\ge2M_0H\), an interval since the exit radius
decreases with \(h\). Fix \(M_1\ge2M_0\) as required by the corollary; this gives the
large-exit interval \(\varrho_e\ge M_1H\).

\subsubsection{Covering the central heights}

The bounded-radius estimate allows any fixed \(M\), whereas the
large-exit threshold \(M_1\) uses only the independently fixed negative
endpoint bound. We may therefore choose \(M\) sufficiently large after
\(M_1\). We show that every height outside the bounded-radius range
satisfies the large-exit estimate. If the
positive half-orbit of a central height enters the adaptive range
thin and has not exited when its radius reaches \(MH\), its
eventual exit is larger.
If it has exited, its radius before cone landing is at most
\(2\varrho_e\). After landing the cone estimates give, whenever the
fixed output height is reached,
\[
 U_+(H^2Y_f)\le C\varrho_e^2+C H^2Y_f .
\]
If the half-orbit was already above the thin threshold at \(M_0H\) but
below the cone, \cref{lem:direct-adaptive-positive} bounds its landing
radius by \(2M_0H\), and the same inequality holds with
\(\varrho_e=M_0H\). If already in the cone, its continuation gives this
bound directly.
Thus, for a sufficiently large fixed \(M\), every central height
outside the bounded-radius range must have \(\varrho_e\ge M_1H\).

The early interval \eqref{eq:direct-primary-early}, the bounded-radius
interval, and the large-exit interval cover \(\cH_p\), with fixed
overlaps. On these intervals the same function \(\mathcal L\) of \(T\)
satisfies \(\mathcal L<0\) in the early and large-exit ranges, and
\(\mathcal L''<0\) on the bounded-radius interval, including the
shared cuts. \Cref{cor:ordered-cover-zero-count} gives at most three
zeros of \(\Delta_p\). This proves \cref{prop:direct-primary-negative}.
